% =========================================================================
\section{A three-mean support of order nine}\label{sec:mean}

\keyidea{Geometric idea}{%
The decentered arc matches one coefficient of the exact target $q_*/\sqrt{t}$ using one radical and one
arc. Combining three members of the same family---the arithmetic mean $\mathrm A=(1+t)/2$, the geometric
mean $\mathrm G=\sqrt{t}$ and one conic mean $\mathrm C_d=\sqrt{t+d(t-1)^{2}}$---gives three free
parameters and matches \emph{three} coefficients, jumping straight to contact eight and order nine.
Geometrically, the transport line $ME$ delivers the residual to three prepared verticals for free, so the
entire support costs only \emph{three compass arcs and one join}---just two traces more than Halley or
the decentered arc. Analytically, changing variable to $\tau=\sqrt{t}$ leaves only one square root, so
global convergence is proved by a single rationalization.}

\subsection{Formula}

The exact support is $\kappa_*=1+(p-1)q_*$ with
\[ q_*(t)=\frac{t-t^{h}}{(p-1)(t^{h}-1)},\qquad q_*(1)=1,\qquad q_*(1/t)=q_*(t)/t , \]
a homogeneous symmetric mean of $1$ and $t$; for integer $p$ it is the arithmetic mean of the $p-1$
mean proportionals $t^{k/p}$ between $1$ and $t$. For real $p>2$ define
\begin{equation}\label{eq:mcoeff}
\begin{gathered}
m=\frac{(p-2)(47p^{2}+21p+4)}{60\,p\,(4p^{2}-1)},\qquad
g=\frac{16\,(p+1)(4p^{2}-1)}{15\,p\,(5p^{2}+4)},\\[2pt]
n=\frac{49\,p\,(p+1)(p^{2}-4)}{10\,(4p^{2}-1)(5p^{2}+4)},\qquad
d=\frac{5p^{2}+4}{84\,p^{2}} ,
\end{gathered}
\end{equation}
all positive, with $2m+g-n=1$, and
\begin{equation}\label{eq:qM}
q_M(t)=m\,(1+t)+g\sqrt t-n\sqrt{t+d\,(t-1)^{2}},\qquad
\Tm_{M}(s)=s\,\frac{1+(p-1)q_M(t)}{t+(p-1)q_M(t)},
\end{equation}
that is, the support $\kappa=1+(p-1)q_M(t)$ in \eqref{eq:support}. With the arithmetic mean
$\mathrm A=(1+t)/2$, the geometric mean $\mathrm G=\sqrt t$ and the \emph{conic mean}
$\mathrm C_d=\sqrt{\mathrm G^{2}+4d(\mathrm A^{2}-\mathrm G^{2})}$,
\[ q_M=2m\,\mathrm A+g\,\mathrm G-n\,\mathrm C_d , \]
a combination with weights of sum one; one weight is negative. The radicand $t+d(t-1)^{2}$ is positive
for every $t>0$, so $q_M$ is defined on the whole positive axis, and $q_M(1/t)=q_M(t)/t$: the log-error
map is odd. In terms of $h=1/p$,
$m=\frac{(1-2h)(4h^{2}+21h+47)}{60(4-h^{2})}$, $g=\frac{16(1+h)(4-h^{2})}{15(4h^{2}+5)}$,
$n=\frac{49(1+h)(1-4h^{2})}{10(4-h^{2})(4h^{2}+5)}$ and $d=\frac{4h^{2}+5}{84}$. For cube roots,
\begin{equation}\label{eq:qMcubic}
q_M(t)=\frac{7}{90}(1+t)+\frac{64}{63}\sqrt t-\frac{6}{35}\sqrt{t+\frac{7}{108}(t-1)^{2}} .
\end{equation}

The three terms belong to one family. The decentered arc is $q_5^{2}=t+a_5(t-1)^{2}$ with
$a_5=(p-2)/(12p)$, as \eqref{eq:arc} shows after expanding $S(t)=(p-2)(t-1)^{2}+12pt$; the members
$a=\tfrac14$ and $a=0$ of the family $q^{2}=t+a(t-1)^{2}$ are $\mathrm A$ and $\mathrm G$, and
$\mathrm C_d$ is the member $a=d$. Each member divided by $\sqrt t$ is $\sqrt{1+aw}$, a function of
$w=(t-1)^{2}/t$ alone, as is $q_*/\sqrt t$. A combination of members therefore matches $q_*/\sqrt t$
coefficient by coefficient in $w$: the arc, with one free parameter, matches one coefficient, and
\eqref{eq:qM}, with three free parameters (four coefficients, one normalization), matches three. The two
members $\mathrm A$ and $\mathrm G$ are imposed; only the third is chosen by the matching.

\subsection{Construction by three arcs}\label{sec:mean-construction}

Let $\lambda_X=H/X$, so that $M=(0,H-\lambda_X)$ and $E=(W,H-\lambda_X t)$. Put
\[ \gamma=1+2(p-1)m,\qquad \zeta=\frac1{2d}-1>1,\qquad K_g=\frac{(p-1)g}{\gamma}\,\lambda_X,\qquad
K_n=\frac{(p-1)n\sqrt d}{\gamma}\,\lambda_X , \]
and prepare four verticals, the horizontal through $M$ and three points:
\[ \begin{gathered}
v:\ x=W\Bigl(1-\frac1\gamma\Bigr),\qquad v_1:\ x=x_1=\frac{(p-1)m}{\gamma}W,\qquad
v_2:\ x=x_2=\frac{K_g}{2\lambda_X}W,\qquad v_3:\ x=x_3=\frac{K_n}{\lambda_X}W,\\[2pt]
F=(x_2,\,H-\lambda_X),\qquad G_2=(x_2,\,H-\lambda_X+K_g),\qquad G_3=(x_3,\,H-\lambda_X+K_n(1+\zeta)).
\end{gathered} \]
The vertical $v$ is at abscissa $2x_1$. The construction below requires that the distance from $v_1$ to
$v$ equal $K_n\sqrt{\zeta^{2}-1}$, which fixes the shape of the rectangle:
\begin{equation}\label{eq:aspect}
\frac{WX}{H}=\frac nm\sqrt{\frac{1-4d}{4d}} .
\end{equation}
The right side is a constructible number depending on $p$ only; it equals $3.7256\ldots$ for $p=3$ and
decreases to $2.238\ldots$ as $p\to\infty$.

Figure~\ref{fig:mean} walks through the four mobile stages in sequence. The line $ME$ is drawn for the
transport; its intersections $J_i=ME\cap v_i$ are free, and their heights are affine in $t$:
$y_{J_i}=H-\lambda_X+(x_i/W)\lambda_X(1-t)$ (panel~1). Three arcs follow:
\begin{enumerate}[label=\arabic*.]
\item \emph{Conic radical} (panel~2). With centre $J_1$ and compass opening $J_3G_3=K_n(t+\zeta)$, draw an
  arc and take its upper intersection $D_1$ with $v$. Its height above $J_1$ is
  $K_n\sqrt{(t+\zeta)^{2}-(\zeta^{2}-1)} = K_n\sqrt{(t+d(t-1)^{2})/d}$. This is the decentered-arc device,
  with a moving centre.
\item \emph{Mean proportional} (panel~3). With centre $J_2$, draw the arc through $G_2$; its radius is
  $\tfrac12K_g(1+t)$ and $J_2$ is at distance $\tfrac12K_g|1-t|$ from the horizontal through $M$, so the
  arc meets that horizontal at a point $Z$ with $FZ=K_g\sqrt t$.
\item \emph{Transfer and readout} (panel~4). With centre $D_1$ and opening $FZ$, draw an arc and take its
  lower intersection $D_M$ with $v$.
\end{enumerate}
Then
\[ H-y_{D_M}=\frac{\lambda_X}\gamma\bigl[1+(p-1)m(1+t)\bigr]-K_n\sqrt{\frac{t+d(t-1)^{2}}{d}}+K_g\sqrt t
=\frac{\lambda_X}\gamma\bigl[1+(p-1)q_M(t)\bigr], \]
and since $W-x_{D_M}=W/\gamma$ the slope of $AD_M$ is $\tfrac{H}{WX}[1+(p-1)q_M(t)]$: the join $AD_M$
is the support. Both radicands are positive for $t>0$, so the three intersections exist at every state.
The mobile work is three arcs and one join, against one arc and one join for the decentered arc. Without
the condition \eqref{eq:aspect} the centre of the first arc is obtained by one additional join, projecting
$E$ from a prepared point onto the vertical at distance $K_n\sqrt{\zeta^{2}-1}$ from $v$; we count the
protocol with \eqref{eq:aspect}.

\begin{figure}[htbp]
\centering\figMeanSteps
\caption{Step-by-step construction of the three-mean support for $p=3$, $X=2$, $H=4$, $s=3/4$, with
$W=7.451\ldots$ from \eqref{eq:aspect}, drawn from exact coordinates at equal axis scales. In each panel
only the active step is highlighted; prior elements recede in grey.
\textbf{(1)}~The transport line $ME$ intersects the prepared verticals $v_1,v_2,v_3$ at the free points
$J_1,J_2,J_3$.
\textbf{(2)}~Arc~1 (centre $J_1$, opening $J_3G_3$) strikes $v$ at $D_1$ via the shaded right triangle.
\textbf{(3)}~Arc~2 (centre $J_2$, through $G_2$) strikes the horizontal of $M$ at $Z$, producing the
geometric-mean segment $FZ=K_g\sqrt{t}$.
\textbf{(4)}~Arc~3 transfers $FZ$ downward from $D_1$ to $D_M$ on $v$; the blue join $AD_M$ is the support
and the green parallel reads $P_+$.}\label{fig:mean}
\end{figure}

\subsection{Global convergence and order nine}

\begin{theorem}[Three-mean support; formalized in Lean]\label{thm:mean}
For real $p>2$, $X>0$ and every $s_0>0$:
\begin{enumerate}[label=\textup{(\roman*)}]
\item $q_M(t)>0$ for every $t>0$;
\item the iteration $s_{j+1}=\Tm_M(s_j)$ converges to $\alpha$ monotonically without crossing it;
\item the local order is exactly nine,
\begin{equation}\label{eq:cM}
e_+=c_M(p)\,e^{9}+O(e^{11}),\qquad
c_M(p)=\frac{(p^{2}-1)(p^{2}-4)(4p^{2}-1)(17p^{2}+8)}{2^{11}\,3^{4}\,5^{2}\,7^{2}}>0 ,
\end{equation}
and $c_M(3)=23/20736$.
\end{enumerate}
\end{theorem}

The constant grows like $p^{8}$, which is the natural scale: in the residual logarithm $x=\log t=pe$ the
law reads $x_+=c_M(p)p^{-8}x^{9}+\dots$, and $c_M(p)p^{-8}\to17\cdot2^{-9}3^{-4}5^{-2}7^{-2}$, about
$3.4\times10^{-7}$, as $p\to\infty$. The same rescaling of \eqref{eq:arc-order} gives the limit $1/480$
for the decentered arc. The asymptotic regime is $|x|\lesssim1$, that is $|e|\lesssim1/p$, for both.

\begin{proof}
Work in the variable $\tau=\sqrt t>0$, in which only one radical remains. Put $c=p-1$ and
\[ S_M(\tau)=\tau^{2}+d\,(\tau^{2}-1)^{2},\qquad \rho_M=\sqrt{S_M},\qquad
\ell_M(\tau)=m(1+\tau^{2})+g\,\tau,\qquad q(\tau)=\ell_M(\tau)-n\,\rho_M(\tau), \]
so that $q_M(t)=q(\sqrt t)$ and $S_M\ge\tau^{2}>0$.

\emph{Positivity.} Clearing denominators gives the identity
\begin{equation}\label{eq:mean-pos}
\ell_M^{2}-n^{2}S_M=\frac{Q_0\,(\tau-1)^{4}+Q_1\,\tau\,(\tau-1)^{2}+Q_2\,\tau^{2}}{900\,p^{2}(4p^{2}-1)(5p^{2}+4)},
\end{equation}
with $Q_0=(p-2)^{2}(2p-1)(313p^{3}+272p^{2}+200p+16)$,
$Q_1=60\,p\,(p-2)(2p-1)(71p^{3}+74p^{2}+40p-8)$ and
$Q_2=180\,p^{2}(149p^{4}+49p^{3}-141p^{2}-196p-20)$. For $p>2$, $Q_0\ge0$, $Q_1>0$ and $Q_2>0$, the
last because
$149p^{4}+49p^{3}-141p^{2}-196p-20=149\varrho^{4}+1241\varrho^{3}+3729\varrho^{2}+4596\varrho+1800$ with
$\varrho=p-2$. Hence $\ell_M^{2}>n^{2}S_M$, and as $\ell_M>0$ we get $\ell_M>n\rho_M$: this is (i).

\emph{The slope.} Let
\[ \Lambda_M(\tau)=\log\bigl(\tau^{2}+c\,q\bigr)-\log\bigl(1+c\,q\bigr)-\frac2p\log\tau ,
\qquad\text{so that}\qquad -\Lambda_M(\sqrt t)=\frac{\log t}{p}+\log\frac{1+c\,q_M}{t+c\,q_M}=e_+ . \]
Both logarithms are defined by (i). Differentiating, with $q'=\ell_M'-nS_M'/(2\rho_M)$,
\[ \tau\rho_M(\tau^{2}+cq)(1+cq)\,\Lambda_M'
=2\tau^{2}(1+cq)\rho_M+c\,\tau(1-\tau^{2})\bigl(\ell_M'\rho_M-\tfrac12nS_M'\bigr)
-\tfrac2p(\tau^{2}+cq)(1+cq)\rho_M . \]
The right side is a cubic polynomial in $\rho_M$. Reducing it with $\rho_M^{2}=S_M$ and substituting
\eqref{eq:mcoeff} gives
\begin{equation}\label{eq:mean-slope}
\Lambda_M'(\tau)=-\,\frac{(p^{2}-4)(p^{2}-1)\,\bigl(2\Psi_2\,\rho_M-7\Psi_1\bigr)}
{1800\,p^{3}(4p^{2}-1)^{2}(5p^{2}+4)^{2}\;\tau\,\rho_M\,(\tau^{2}+cq)(1+cq)} ,
\end{equation}
where $\Psi_1$ is a palindromic sextic and $\Psi_2$ a palindromic quartic in $\tau$ with coefficients in
$\mathbb Z[p^{2}]$, listed in full in Appendix~\ref{app:certs} (Table~\ref{tab:mean-cert}). Every factor
of the denominator is positive.

\emph{The sign.} The three coefficients of $\Psi_2$ in Table~\ref{tab:mean-cert} are positive for $p>2$,
so $\Psi_2>0$ for $\tau>0$. Rationalizing,
\begin{equation}\label{eq:mean-square}
4\Psi_2^{2}\,S_M-49\Psi_1^{2}=\frac{4\,(4p^{2}-1)^{3}(5p^{2}+4)^{2}}{21\,p^{2}}\;(\tau-1)^{8}\;\Psi_3 ,
\end{equation}
a polynomial identity after clearing $84p^{2}$, where $\Psi_3$ is a third palindromic quartic whose
three coefficients (Table~\ref{tab:mean-cert}) are again positive for $p>2$. Let $\tau>0$, $\tau\ne1$. If
$\Psi_1\le0$ then $2\Psi_2\rho_M>0\ge7\Psi_1$. Otherwise both sides are positive and \eqref{eq:mean-square}
gives $(2\Psi_2\rho_M)^{2}>(7\Psi_1)^{2}$. In both cases $2\Psi_2\rho_M>7\Psi_1$, and by
\eqref{eq:mean-slope} $\Lambda_M'<0$ for $\tau\neq1$. The factor $(\tau-1)^{8}$ is the contact of order
eight.

\emph{Convergence.} $\Lambda_M$ is strictly decreasing and $\Lambda_M(1)=0$, because
$q(1)=2m+g-n=1$. Hence $e_+=-\Lambda_M(\sqrt t)$ has the sign of $t-1$, that is the sign of $e$: the
state does not cross $\alpha$. Meanwhile $\Tm_M(s)/s-1=(1-t)/(t+c\,q_M)$ moves the state toward
$\alpha$. The orbit is monotone and bounded between $s_0$ and $\alpha$; its positive limit is a fixed
point by continuity, and a fixed point has $t=1$. This is (ii).

\emph{Local order.} Both $q_M/\sqrt t$ and $q_*/\sqrt t$ are even analytic functions of $\log t$, hence
power series in $w$, and exact expansion gives
\begin{equation}\label{eq:mean-D4}
\frac{q_M-q_*}{\sqrt t}=D_4^{M}\,w^{4}+O(w^{5}),\qquad
D_4^{M}=\frac{(2-h)(1+h)(2+h)(1-2h)(1+2h)(8h^{2}+17)}{203212800} .
\end{equation}
Since $w=\xi^{2}/t$, $\kappa-\kappa_*=(p-1)D_4^{M}\xi^{8}+O(\xi^{9})$, and Lemma~\ref{lem:contact} with
$\omega=8$ gives $e_+=(p-1)p^{7}D_4^{M}e^{9}+O(e^{10})$, which is \eqref{eq:cM}; oddness removes the
term in $e^{10}$. In Lean the expansion \eqref{eq:mean-D4} is obtained without Taylor series. Put
$y=\tfrac12\log t$, so that $\sqrt t=e^{y}$ and the conic radical is $e^{y}\sqrt{1+dw}$. If $R_4$ is the
quartic Taylor polynomial of $\sqrt{1+dw}$, then $1+dw-R_4^{2}$ is $w^{5}$ times a polynomial, so
replacing the radical by $e^{y}R_4(w)$ costs $O(w^{5})$. After this replacement and after subtracting
$D_4^{M}e^{y}w^{4}$, the difference with $q_*$, multiplied by
$60(4-h^{2})(4h^{2}+5)(1-h)(e^{2hy}-1)e^{8y}$, is a sum of $22$ exponentials $c_j(h)e^{(k_j+2hl_j)y}$
whose moments $\sum_jc_j(k_j+2hl_j)^{n}$ vanish for $n<11$; it is therefore $O(y^{11})$, and
$w/y^{2}\to4$ gives $(q_M-q_*)/y^{8}\to256\,D_4^{M}$.
\end{proof}

The proof has the same structure as the proof of Theorem~\ref{thm:arc}---one radical, one rationalization
and one sign discussion: although $q_M$ contains two radicals, in the variable $\sqrt t$ only one remains,
and three quartics suffice.

\begin{remark}[Other members of the family; computer algebra]\label{rem:family}
The same matching with fewer members gives lower orders. With $\mathrm A$ and one free conic mean,
\[ q_7(t)=m_7\,(1+t)+n_7\sqrt{t-\varepsilon_7\,(t-1)^{2}},\qquad
m_7=\frac{(h+3)(1-2h)}{4(4-h^{2})},\quad n_7=\frac{5(1+h)}{2(4-h^{2})},\quad
\varepsilon_7=\frac{1-4h^{2}}{60}, \]
has contact six: a single radical gives order seven, with
$e_+=-\tfrac{(p^{2}-1)(p^{2}-4)(4p^{2}-1)}{302400}e^{7}+O(e^{9})$; for cube roots
$q_7=\tfrac1{14}(1+t)+\tfrac67\sqrt{t-(t-1)^{2}/108}$. But the conic parameter is negative: the radicand
is positive only for $w<60/(1-4h^{2})$, and the negative constant means that the state crosses the root
at each step. The same happens for the combination of two free conic means, which has contact eight
with one negative parameter. Imposing both $\mathrm A$ and $\mathrm G$ is what makes the third
parameter $d$ positive in \eqref{eq:qM}. These statements are exact series computations; no
convergence theorem is claimed for $q_7$.
\end{remark}
